\chapter{Introduction}

\section{What this report is}

This report explains a complete software project from end to end: what problem
it solves, how it was designed, how each part works, how it is deployed, and how
it is kept running once nobody is watching it. It is written as a guide, so a
reader who has never seen the code can follow the reasoning and learn from it.

The project is a \textbf{personal finance manager}: a small web application that
records expenses and income, tracks investments across several brokers,
calculates net worth, produces Excel and PDF reports, and answers from a
Telegram chat.

Every chapter follows the same shape. It states the problem, shows the decision
that was taken, and explains why the alternatives were left aside. Code appears
only when it makes the explanation shorter than prose would.

\section{Who it is for}

The report assumes the reader knows Python and has seen a web framework before.
It does not assume knowledge of FastAPI, SQLAlchemy, Alembic, the Telegram Bot
API or financial concepts. Those are introduced where they first appear.

A reader who only wants to run the application does not need this document; the
repository \file{README.md} covers that in a page. This report is for the reader
who wants to understand \emph{why} the code looks the way it does.

\section{Where the project came from}

The application replaced five desktop scripts written in Python with a
\code{tkinter} interface. Each script covered one area: expenses, income,
investments, a shared property, and account balances. They worked, but they
carried the problems that small personal tools usually carry:

\begin{itemize}
  \item They only ran with the computer switched on, so nothing could be
        recorded from a phone.
  \item Data lived in loose files, mostly spreadsheets, text files and PDFs.
        There was no database, no history, and no safe way to query anything.
  \item Amounts were handled as \code{float}. Cents drifted.
  \item Personal data and absolute paths were written directly in the source.
  \item There were no tests, so a change in one script could break another
        without anybody noticing.
\end{itemize}

The goal of the rewrite was a single application, reachable from a phone, with a
real database, sound arithmetic, and no personal data in the code.

\section{What the application does today}

\begin{table}[h]
\centering
\begin{tabular}{@{}p{4.1cm}p{9.6cm}@{}}
\toprule
\textbf{Area} & \textbf{What it covers} \\
\midrule
Expenses and income & Manual entry plus a paste importer that reads the notes
the user already keeps on a phone, evaluates the arithmetic in each line, and
shows a preview before saving. \\
Investments & Statement import for three brokers, manual positions, automatic
revaluation from market prices, and recurring contributions. \\
Net worth & Bank accounts and cash added to the investment value, with a daily
snapshot that feeds an evolution chart. \\
Reports & Excel exports and PDF reports with charts, for a month, a year or any
range of months. \\
Telegram bot & Logging entries from a chat with a confirmation step, plus
summaries, balances and reports on demand. \\
Operations & Health endpoints, scheduled jobs, and a weekly database backup
delivered outside the hosting provider. \\
\bottomrule
\end{tabular}
\caption{Scope of the application.}
\end{table}

\section{Design principles}

Five rules shaped almost every decision in the project. They are stated here
once and referenced throughout the report.

\begin{keypoint}[The five rules]
\begin{enumerate}
  \item \textbf{Money is never a float.} Amounts are \code{Decimal} in Python
        and \code{NUMERIC} in the database.
  \item \textbf{Nothing personal in the code.} Every value that changes between
        users comes from the environment.
  \item \textbf{One place per calculation.} The web pages, the reports and the
        bot call the same functions, so their totals cannot disagree.
  \item \textbf{Fail loudly where it matters, quietly where it does not.} A
        failed backup returns an error; a failed price lookup leaves the old
        value in place.
  \item \textbf{Boot fast.} Slow startup is not a comfort issue; on a host that
        sleeps, it is the difference between a bot that answers and one that
        does not.
\end{enumerate}
\end{keypoint}

\section{How to read this report}

Chapters \ref{ch:architecture} to \ref{ch:reports} describe the application from
the inside: structure, data, and each functional area. Chapter
\ref{ch:bot} covers the Telegram bot. Chapters \ref{ch:operations} to
\ref{ch:deployment} cover everything that happens after the code is written:
running it, securing it, testing it and deploying it. Chapter
\ref{ch:going-public} describes a separate exercise, turning a private
repository into a public one without leaking anything. The last chapter collects
the lessons that were only obvious afterwards.
